\section{The Multiverse and Massachusetts}

I had the opportunity this weekend to watch an interview on C-Span with the physicist Brian Greene\footnote{\url{http://www.c-spanvideo.org/videoLibrary/event.php?id=192289\&timeline}}, conducted at the Museum of Science in Boston, Massachusetts. The topic was the multi-worlds interpretation of quantum mechanic and the all-white audience was mostly students and retired folk. I know that tribe intimately, having grown up and attended University in the Boston area. I'll touch on the science of the theory, but what is more interesting is the sociology of the audience.

\paragraph{Theory}


The theory itself is quite simple: whenever there is quantum indeterminancy, then both possibilities occur, one in ``our'' universe, and the other in an alternate universe that is identical except that the other possibility is observed. I can understand why Mr. Greene is excited about this theory. Had he given a lecture on his specialized field of string theory, perhaps a dozen people would have been able to understand it and it would not have been televised. But when it comes to the multiverse, suddenly everyone is an expert on quantum mechanics. When scientists start approaching metaphysical topics, I don't know if they are simply over their heads, or if they just get some satisfaction in deluding the audience.

The theory itself depends on a mathematical trick, viz., if an equation has multiple solutions, then they all must refer to physical reality. Yet, that is not the case in other branches of physics. For example, in systems that involve the solution to a quadratic equation, the negative solution is rejected as not conforming to physical reality, while the positive solution is accepted. The multiverse theory asserts that the unobserved (and unobservable) solution also physically exists. That remains to be seen, since physicists qua physicists must also experimentally demonstrate their theories, not just theoretically.

Mr. Greene curiously stated an obvious error in physics. The question arose about the fundamental constants in physics\footnote{\url{http://en.wikipedia.org/wiki/Physical_constant}} whose values are really inexplicable to current physics. Mr. Greene then postulated that there are alternate universes with different fundamental constants. However, the multiverse theory depends on quantum effects and no one has demonstrated, or even theorized, that the fundamental constants are in any way related to quantum effects. Hence, there is no reason to accept that possibility, even as conjecture.

He then made an allusion to the Anthropic Principle\footnote{\url{http://en.wikipedia.org/wiki/Anthropic_principle}}, with the unsupported claim that these other universes would not have been such that life could have evolved, and ``we'' happen to be in the universe with the right constants. Of course, an unobservable universe is no longer physics; it is metaphysics, and not even good metaphysics. Mind and matter are correlative; as Jung put it in \emph{Answer to Job}, ``Existence is only real when it is conscious to someone.''

\paragraph{Sociology}


Even if a physicists can create alternative universes at will by sending electrons through a slit, that is not necessarily of great significance. So suddenly the theory morphs into the idea that whenever a choice is made, then another universe is created where the opposite choice was made. For this to be the case, we would have to assume that the human organism as a whole is a quantum effect, and that ``making a choice'' is akin to the seemingly random motion of an electron beam through a narrow slit. Did Mr. Green propose any serious scientific reason to believe that? But as Chesterton said, ``if a man does not believe in God, he will believe an anything.'' Certainly the audience seemed to believe in anything.

Just as believers in reincarnation assume they were great leaders or princesses in their past live, believers in the multiverse presume ``they'' will have made great achievements in an alternative universe. One questioner seemed to think that in some other universe, he is a great physicist rather than an engineer on Route 128. The question of his real identity did not seem to cross his mind. Of course, in that audience, Mr. Greene made the obvious joke, to great laughter, about Sarah Palin being president in some other universe. Not that I have any particular affection or respect for Mrs. Palin, the aim of the joke was to establish a tribal identity, not to encourage independent thought.

Had the liberal members of the audience taken a pause to think things through, they should have come to the conclusion that in alternative universes, they may just as well be tea party activists, or even neo-Nazi racists. If every conceivable possibility is actual, then it eliminates the necessity for choice and makes a mockery of the heroic human life.


\flrightit{Posted on 2011-04-03 by Cologero}

\begin{center}* * *\end{center}

\begin{footnotesize}
\begin{sffamily}
\texttt{Jake on 2011-05-03 at 20:28 said:}

You might want to check out some of the more bleeding-edge stuff going on with MWI -- there are a few splinters who are looking at the big picture in a philosophically coherent way. One of the more exciting approaches I've seen points out that under mathematical information theory, the set of all possible descriptions -- strings of 1s and 0s, essentially speaking -- has a total information content of precisely zero, and from there proposes that Something (the manifestation of a determinate reality in temporal appearance) emerges from the primordial All/Nothing, now formally understood as the set of all possible worlds, by means of conscious observation. Obviously, the sense and usefulness of this is crippled by the need for a theory of the observing, reflexive subject(ive process) -- this is nowhere more evident than in the book's discussion of time -- but I think there's a lot that could be done with the synthesis of MWI, information theory, and modern continental ontology, particularly that of Hegel, Heidegger, and Deleuze.

\hfill

\texttt{Cologero on 2011-05-03 at 22:28 said:}

You seem to be referring to a book, but the title is omitted. Does it address the following issues:

1) Is the set of all possible descriptions finite or uncountable?

2) Is any particular universe bounded by a maximal string size?

3) Is every permutation a possibility of manifestation, or are some permutations simply impossible?

4) Or, more likely, are some strings ``true'' but not ``provable''?

5) If a string represents a state of the universe, what are the rules for moving from state to state? In other words, is the set of all possible descriptions a group or algebra?

6) ``0'' is the initial state, and it transitions to ``1''. Then the ``1'' to the many. Correct? This is the basis of every metaphysical system. How does MWI improve on this?

7) Finally, metaphysics has its own version of many worlds. It is also true that all possible worlds will manifest. Physics doesn't make it so, actually, just the opposite.

\hfill

\texttt{Jake on 2011-05-04 at 07:54 said:}

Sorry about the lack of title -- it's THEORY OF NOTHING, written by Russell Standish. It's available online:
\url{http://www.hpcoders.com.au/nothing.html}


1) No -- the set of all possible descriptions is uncountable from definition, I would think.

2) I think this probably depends on how we mean `universe' -- there's at least two important and distinct modes. If we mean by it the temporally organized world manifest by a particular observer, then yes, all world-narratives should eventually terminate or decohere. If we mean `universe' in the sense of the scientifically ascertained cosmos, then I'd say it's a pretty open question.

3) Well, your first test is clearly anthropic -- it has to (be able to) describe observer processes in order to become actual -- but obviously no one has any real clue what that entails. Beyond that, Standish uses statistics and algorithmic information theory at this point, to try to demonstrate that descriptions corresponding to simpler rulesets are more likely to be observed and ``prove'' the Razor. It's a cool idea, but the ontological ambiguity introduced with rulesets, algorithms, and correspondences makes me very skeptical.

4) Not sure how truth and provability in the metamathematical senses can be applied here... can you elaborate?

5) Personally, I would refer this to the need for a theory of observation -- but I have certain Heideggerian notions concerning the relation of observation and temporality which you may not share.

6) There's no transition. Transitions are temporal, and the relation of the particular observed world to prior Totality (aka All/Nothing) is pre-temporal. Since we identify Nothing with indeterminate Totality rather than a state of absolute absence, the ``why'' of Something is no longer problematic; it is intrinsically potential within Nothing, and is necessarily manifest in actuality insofar as it describes an observer.

7) Not really sure what you mean here -- Leibniz? I would argue that metaphysics is the only discipline which is even capable of talking about worlds. Most MWI is just bad metaphysics that's not even aware of itself as such.

\end{sffamily}
\end{footnotesize}

